\documentclass[11pt]{article}
\usepackage[T1]{fontenc}
\usepackage[margin=1in]{geometry}
\usepackage{amsmath,amssymb,amsthm,booktabs,longtable,array}
\usepackage[colorlinks=true,allcolors=blue]{hyperref}
\usepackage{microtype}
\newtheorem{definition}{Definition}[section]
\newtheorem{proposition}[definition]{Proposition}
\newtheorem{remark}[definition]{Remark}
\newcommand{\R}{\mathbb R}
\newcommand{\C}{\mathbb C}
\newcommand{\Tr}{\operatorname{Tr}}
\newcommand{\id}{\mathbb I}
\newcommand{\Lset}{\mathcal L}
\newcommand{\Q}{\mathcal Q}
\newcommand{\NS}{\mathcal{NS}}
\newcommand{\val}{\omega}
\newcommand{\ket}[1]{\lvert #1\rangle}
\newcommand{\bra}[1]{\langle #1\rvert}
\title{NonLocaLang: Resource-Aware Primitives for Nonlocal Games}
\author{NonLocaLang research project}
\date{}
\begin{document}
\maketitle
\begin{abstract}
A useful abstraction for nonlocal games must preserve the players' information and
the admissible resources, not only the referee's score. We propose a layered
representation separating the experimental arena, verifier, resource model,
strategy witness, and certification question. Finite classical one-round games
have an exact tabular representation and several useful algebraic constructors.
Independent sources, trusted quantum interfaces, and sequential interactions
require additional semantic structure. We establish elementary preservation and
obstruction results that constrain sound compiler passes, organize numerical and
symbolic analyses by the direction of their guarantees, and specify a concrete
implementation scope. A central distinction is between certifying a value,
certifying a behavior, and extracting a target state or measurement. The framework
is a research architecture with a working finite-game core, not a universal
decision procedure or a claim that arbitrary games have efficient analyses.
\end{abstract}
\begingroup
\small
\setcounter{tocdepth}{1}
\tableofcontents
\endgroup
\clearpage
\section{What a game specification must preserve}
\label{sec:semantics}

The standard nonlocal game is a cooperative test: a referee samples questions,
sends each player only their own question, receives answers, and evaluates a
score. Its meaning depends on the absence of communication during the round and
on what resources were distributed beforehand \cite{cleve2004}. Entanglement is
a resource that can generate Bell-nonlocal correlations; entanglement and Bell
nonlocality are not interchangeable properties. CHSH violation certifies
entanglement under its assumptions, but failure to violate a particular Bell
inequality does not certify separability.

\begin{definition}[Finite classical one-round verifier]
An arena is a finite list of parties $i\in\{1,\ldots,n\}$, each with a nonempty
question alphabet $X_i$ and answer alphabet $A_i$. Write
$X=\prod_iX_i$ and $A=\prod_iA_i$. A verifier is a distribution
$\mu:X\to[0,1]$ and a real payoff $V:A\times X\to\R$. A winning-probability
verifier has $0\leq V\leq1$; a deterministic win/lose predicate takes values in
$\{0,1\}$. A behavior is a normalized conditional distribution $p(a\mid x)$,
and its score is
\begin{equation}
 S_G(p)=\sum_{x\in X}\mu(x)\sum_{a\in A}V(a,x)p(a\mid x).
 \label{eq:score}
\end{equation}
The standard information rule gives player $i$ only $x_i$, with no interparty
communication during the round. Question sampling is independent of the
players' pre-existing resources.
\end{definition}

There is no invariant answer to ``how many fields define a game'': one may
encode all data in a single string, infer $n$ from an alphabet list, or combine
$\mu$ and $V$. A practical answer has five semantic components:
\begin{enumerate}
\item an arena with parties, interfaces, and their information/causal structure;
\item a verifier containing sampling, acceptance/payoff, and any trusted operations;
\item an admissible resource/strategy model $\mathcal R$;
\item an analysis question, such as a value, membership, separation, or extraction bound;
\item optional reference realizations and a target-equivalence specification.
\end{enumerate}
The first two specify the test. Adding the third defines the optimization
problem $\val_{\mathcal R}(G)=\sup_{p\in\mathcal R}S_G(p)$. Reference states
are not necessary to define $G$ or its value. They become essential when the
question is whether a score certifies a particular realization. The supremum
must not silently become a maximum.

Question-dependent answer alphabets can be represented using a common padded
alphabet only with explicit legality restrictions. Merely assigning zero payoff
to invalid answers is unsafe for signed payoffs: an invalid answer could become
an advantageous abort. The first implementation uses a fixed answer alphabet
for each party. A party with no external input has $|X_i|=1$; this is a syntax
convention, not permission to erase source-independence assumptions.

\subsection{An exact core representation, with a deliberate boundary}
Define $c_{a,x}=\mu(x)V(a,x)$. For a fixed arena,
$S_G(p)=\langle c,p\rangle$ is an exact intermediate representation (IR) for
linear score evaluation. Every finite classical one-round verifier has this
representation, and a bounded table can be serialized exactly when its entries
are rational. It does not follow that every physically meaningful experiment
is such a verifier, or that arbitrary real parameters have finite exact
encodings.

The implementation should retain both $\mu$ and $V$ alongside $c$. Different
sampling schemes can have the same expected linear score yet different
estimator ranges, variances, and experimental costs. Constraints on the
resource set also refer to the arena rather than just to coefficients. In
particular, a question tuple with zero sampling probability still belongs to
the arena; a non-signalling extension may impose constraints involving it.
Dropping those questions is a separate transformation requiring justification.

There are several inequivalent notions of ``the same game.'' Exact transcript
equivalence preserves the distribution and score of each run. Score
equivalence preserves $S_G(p)$ for every admissible $p$. Value equivalence
preserves only an optimum for one specified resource model. A relaxation
preserves neither behavior nor optimum, but supplies an inequality. A
compiler pass should state which of these contracts it provides.

A coefficient table costs $|X||A|$ scalar entries. A compact parity predicate
can be exponentially smaller. The universal finite table is therefore a
\emph{semantic fallback}, not a zero-cost implementation. Structured
expressions---parity tests, consistency constraints, graph predicates, and
local wirings---should survive until a backend actually needs expansion.

\section{Resources and the limits of a single backend}
\label{sec:resources}

\subsection{Distinct correlation models}
For a fixed finite Bell arena, the unrestricted local set is
\begin{equation}
 p(a\mid x)=\sum_\lambda q(\lambda)\prod_i p_i(a_i\mid x_i,\lambda).
 \label{eq:local}
\end{equation}
Its extreme points are deterministic response functions $a_i=f_i(x_i)$.
Private randomization and shared randomization do not improve a linear optimum
beyond the best deterministic response tuple.

Finite-dimensional tensor-product quantum behaviors have the form
\begin{equation}
 p(a\mid x)=\Tr\left[\rho\bigotimes_i M^{(i)}_{a_i\mid x_i}\right],\qquad
 M^{(i)}_{a_i\mid x_i}\succeq0,\quad
 \sum_{a_i}M^{(i)}_{a_i\mid x_i}=\id.
 \label{eq:quantum}
\end{equation}
The dimension-bounded set $\Q_{\boldsymbol d}$, the union $\Q_{\rm fin}$ over
all finite dimensions, and its closure $\Q_{\rm a}$ are different types.
For a dimension bound, one must also say whether free shared classical
randomness is allowed: convexifying $\Q_{\boldsymbol d}$ can change membership,
although it does not change a linear optimum. Projective-measurement dilation
is harmless when local dimension is unrestricted, but generally changes a
fixed-dimensional resource contract.

In the commuting-operator set $\Q_{\rm c}$, all measurement operators act on a
single Hilbert space, and operators assigned to distinct parties commute.
Non-signalling behaviors satisfy, for every subset $S$ of parties,
\begin{equation}
 \sum_{a_{\bar S}}p(a_S,a_{\bar S}\mid x_S,x_{\bar S})
 \quad\text{is independent of }x_{\bar S}.
 \label{eq:ns}
\end{equation}
The useful inclusion chain is
\begin{equation}
 \Lset\subseteq\Q_{\rm fin}\subseteq\Q_{\rm a}
 \subseteq\Q_{\rm c}\subseteq\NS.
 \label{eq:chain}
\end{equation}
Equalities and strictness depend on the arena. In general $\Q_{\rm fin}$ is not
closed \cite{slofstra2019}, and $\Q_{\rm a}\ne\Q_{\rm c}$ for some finite
arenas \cite{ji2020}. The closure has the same supremum for every continuous
linear score as $\Q_{\rm fin}$, but membership and attainment can differ.
These are operational distinctions, not interchangeable solver labels.

\subsection{Three extension axes}
A useful taxonomy is by additional semantics, not by a fixed number of disjoint
classes. The following axes can occur together.

\paragraph{Independent sources and causal structure.}
A network resource specification contains a source hypergraph. In a classical
network, independent latent variables $\lambda_e$ obey
\begin{equation}
 p(a\mid x)=\int\prod_e dq_e(\lambda_e)
 \prod_i p_i(a_i\mid x_i,(\lambda_e)_{e\ni i}).
 \label{eq:network}
\end{equation}
A quantum network has a product $\bigotimes_e\rho_e$ of source states, with
local measurements on all systems arriving at each node. A global hidden
variable shared across sources destroys the independence assumption. The
resulting sets are generally nonconvex; their analysis requires causal
compatibility methods such as inflation \cite{fritz2012,wolfe2021}.
Forgetting independence is a relaxation, never a lossless normalization.

\paragraph{Trusted quantum interfaces.}
Semiquantum games use known quantum states as questions. Substituting their
classical preparation labels as freely readable inputs grants the players
information they did not possess when the states were nonorthogonal. Thus
finite labels alone do not provide a faithful Bell-game lowering. Buscemi's
semiquantum tests witness entanglement beyond ordinary Bell tests
\cite{buscemi2012}. Quantum answers, steering assemblages, and trusted referee
systems similarly require operator-valued interfaces and explicit trust types.

\paragraph{Multiple rounds and memory.}
Sequential interactions need an event order, local histories, admissible
communication, and quantum instruments or channels. A transcript table alone
does not enforce that an early answer is independent of a future question.
For finite-dimensional quantum interactions, strategy/comb representations
provide positive operators with causal trace constraints
\cite{gutoski2007}. Locality between different players remains an additional
restriction; it is not supplied by positivity of a single global comb.

The proposed umbrella IR is a typed causal process: classical or quantum wires,
trusted or untrusted nodes, independent source factors, and a referee score.
It can describe many finite protocols uniformly, but each backend must accept
only a verified fragment. Continuous interfaces, infinite rounds,
superselection rules, energy constraints, computational restrictions, and
generalized probabilistic theories add further semantics. No claim of a
complete finite classification is made.

\begin{proposition}[No-input linear-score obstruction]
\label{prop:noinput}
Suppose every party has one input and every deterministic joint output is an
admissible strategy in a resource model $\mathcal R$. Then
\[
 \val_{\mathcal R}(V)=\max_{a\in A}V(a).
\]
In particular, ordinary local and quantum models cannot be separated by such
a game. The same value obstruction applies to independent-source networks
when arbitrary constant outputs are permitted.
\end{proposition}
\begin{proof}
Every score is a convex combination of the numbers $V(a)$, hence at most
their maximum. The players attain the maximum by each outputting their
component of one maximizing tuple, without using any source.
\end{proof}
No-input network nonlocality is nevertheless meaningful: it concerns whether
a \emph{distribution} belongs to a nonconvex compatible set
\cite{fritz2012}. Its detection generally needs a nonlinear distribution
witness, constraints fixing other statistics, or several measured quantities.
Such tests must be represented as compatibility tasks rather than misleadingly
reported as an unrestricted scalar winning-probability advantage.

\section{Primitives and proof obligations for composition}
\label{sec:composition}

The useful primitives are not a list of famous inequalities. They are typed
constructors for arenas, samplers, predicates, resources, and transformations.
An XOR predicate, for example, specifies
$\bigoplus_i a_i=f(x)$; a consistency predicate compares overlapping answer
components; a graph predicate checks a relation between vertex labels. CHSH,
GHZ parity tests, and the magic-square game are compositions of these small
rules with distinct arenas and resource witnesses.

\subsection{Lossless transformations}
\begin{proposition}[Local relabeling and affine score changes]
Let each question alphabet and each answer alphabet conditional on the local
question be relabeled bijectively. Transport the sampler, payoff, and strategy
by the same local relabeling. Every resource class closed under these
relabelings has the same game value before and after the transformation.
Moreover, for $\alpha>0$ and an input-only function $b(x)$,
\[
 V'(a,x)=\alpha V(a,x)+b(x)
 \quad\Longrightarrow\quad
 \val_{\mathcal R}(G')=\alpha\val_{\mathcal R}(G)
 +\sum_x\mu(x)b(x).
\]
The maximizing behaviors coincide under the transport.
\end{proposition}
\begin{proof}
A local bijection gives an invertible correspondence between admissible
strategies; changing summation indices preserves every score. Normalization
$\sum_a p(a\mid x)=1$ proves the affine identity pointwise, so it also holds
after taking the supremum.
\end{proof}
Negative scaling swaps maximization and minimization. Scaling a winning
probability does not preserve its statistical interpretation or robustness
parameter: a score deficit $\varepsilon'$ becomes $\varepsilon'/\alpha$.

The distinction between local and global input processing is essential. A
map depending on the complete input tuple may reveal another party's question
and change the resource model. Local stochastic preprocessing/postprocessing
is a safe strategy map for ordinary local and quantum models, but is usually
not invertible and therefore need not preserve values in both directions.

\subsection{Mixtures: whether a tag is visible is part of the game}
For games $G_j$ on a common arena and weights $\lambda_j\geq0$ summing to one,
a \emph{hidden mixture} has coefficient table
$c=\sum_j\lambda_j c_j$. It can retain a verifier by using
\[
 \mu(x)=\sum_j\lambda_j\mu_j(x),\qquad
 V(a,x)=\frac{\sum_j\lambda_j\mu_j(x)V_j(a,x)}{\mu(x)}
\]
when $\mu(x)>0$, with arbitrary payoff at zero-mass questions. The players
receive no extra component tag. The questions themselves may still convey
partial or complete information about the component.

\begin{proposition}[Mixture bound and saturation]
For a hidden mixture $G=\sum_j\lambda_jG_j$,
\begin{equation}
 \val_{\mathcal R}(G)\leq\sum_j\lambda_j\val_{\mathcal R}(G_j).
 \label{eq:mixture}
\end{equation}
If a single strategy simultaneously attains every component optimum with
positive weight, equality holds. If component bounds $u_j$ are valid, and a
strategy has mixture score at least $\sum_j\lambda_j u_j-\varepsilon$, then
\[
 0\leq u_j-S_{G_j}(p)\leq\varepsilon/\lambda_j
 \quad(\lambda_j>0).
\]
\end{proposition}
\begin{proof}
Bound each summand separately. For the last assertion, the weighted sum of
the nonnegative component deficits is at most $\varepsilon$.
\end{proof}
This elementary deficit calculus is useful for rigidity constructions: a
near-optimal combined score forces each consistency test to be nearly
satisfied. It does not establish that the local extraction maps obtained from
different component tests are compatible.

A \emph{public mixture} sends the same component tag $j$ to every party. When
the model permits the necessary conditional strategies and common resource
realization, its value is the weighted sum of the component values (with
suprema interpreted through arbitrarily good strategies). Unrestricted local
and finite-dimensional quantum models permit this by pre-sharing the required
resources and choosing measurements using the tag. Fixed dimension, fixed
state, and independent-source contracts require a separate closure check.
Erasing a public tag or revealing a hidden one is not a compiler optimization.

\subsection{Parallel products and lowerings}
The parallel conjunction $G\otimes H$ sends each player the pair of their
questions, uses independent referee sampling, and accepts iff both component
predicates accept. More generally its payoff is the product of nonnegative
payoffs. If the resource class is closed under independent tensor/product
strategies, then
\[
 \val_{\mathcal R}(G\otimes H)\geq
 \val_{\mathcal R}(G)\val_{\mathcal R}(H).
\]
The product witness proves this inequality. An optimal joint strategy may
couple answers across coordinates; no general multiplicativity equality is
licensed. Parallel conjunction is also different from an XOR combination of
winning conditions, for which specialized theorems may apply.
For example, exact deterministic optimization of two parallel CHSH games
gives classical value $10/16$, which is strictly greater than
$(3/4)^2=9/16$. This benchmark is reproduced by the implementation's
composition tests.

Every proposed lowering should expose three items: an admissible-strategy
map, its score relation, and its resource/size cost. For example, forgetting a
dimension bound or source independence enlarges a feasible set and gives an
upper bound. Restricting to an explicit family of strategies gives a lower
bound. Projective dilation preserves unrestricted one-round quantum values
but consumes dimension. A finite coefficient expansion preserves scores but
may consume exponential memory. These contracts make abstraction cost
visible instead of hiding it behind a uniform function name.

\begin{proposition}[Resource monotonicity]
If $\mathcal R\subseteq\mathcal S$ in the same arena, then
$\val_{\mathcal R}(G)\leq\val_{\mathcal S}(G)$ for every payoff.
Consequently an upper bound for $\mathcal S$ is an upper bound for
$\mathcal R$, and a feasible $\mathcal R$ witness is a feasible
$\mathcal S$ witness.
\end{proposition}
\begin{proof}
The supremum over a subset cannot exceed the supremum over the full set.
\end{proof}

\section{An analysis pipeline with explicit evidence types}
\label{sec:analysis}

\subsection{Values, bounds, and numerical evidence}
An analysis result should contain its resource model, method, lower/upper
direction, precision, feasibility diagnostics, and evidence. A floating-point
solver status of \texttt{optimal} is not a mathematical certificate. In
particular, the primal objective of a numerical maximization is not, by
itself, a rigorous upper bound on the exact problem.

\paragraph{Classical value.}
Deterministic enumeration is exact for rational verifier tables. The number
of deterministic strategy tuples is
$D=\prod_i |A_i|^{|X_i|}$. Enumerating all but one party and computing that
party's optimal answer separately for each question saves a factor without
changing the optimum. This is still exponential in general. Store a maximizing
response table and its exact rational score. A proof of optimality also needs
the exhaustive-search argument or a separately checked dual/branch certificate.

\paragraph{Non-signalling value.}
Maximize $\langle c,p\rangle$ over positivity, normalization, and
Eq.~\eqref{eq:ns}. This finite linear program gives an upper bound for the
models in Eq.~\eqref{eq:chain}. Marginal constraints must include unsupported
question tuples unless their removal has been justified. Numerical feasibility
residuals are informative but do not replace exact dual feasibility. Rational
reconstruction provides a separate certificate path: reconstruct primal and
dual candidates, verify all equalities and inequalities over the rationals,
and compare their objective values. The implementation can certify a rational
optimum when these bounds coincide, or return the certified upper bound
together with the verified interval when they differ. Both primal and dual
feasibility checks must pass. Failed reconstruction is not reported as a proof.

\paragraph{Quantum outer bounds.}
The NPA hierarchy introduces moments
$\Gamma_{u,v}=\langle u^\dagger v\rangle$ of words in measurement operators,
requires $\Gamma\succeq0$, and imposes the algebraic identities of projectors
and interparty commutation. Finite levels are outer relaxations. The full
hierarchy converges to the commuting-operator set, not in general to
$\Q_{\rm a}$ \cite{npa2008}. Thus the same finite-level bound is valid for
finite-dimensional quantum strategies, but its interpretation must state the
limiting model. Restricted words such as $1+A+B$ and $1+A+B+AB$ are distinct
relaxations, not arbitrary-dimension quantum optimizers.
The implemented level-$1$ relaxation also imposes explicit nonnegative
probabilities, so its precise meaning is $Q_1\cap\NS$; this convention matters
when comparing published bounds.

\paragraph{XOR specialization.}
For a two-party binary-output XOR game with winning rule
$a\oplus b=f(x,y)$, define $C_{xy}=\mu(x,y)(-1)^{f(x,y)}$. The quantum winning
probability is $(1+\beta_q)/2$, where
\begin{equation}
 \beta_q=\max_{\|u_x\|=\|v_y\|=1}\sum_{x,y}C_{xy}u_x\cdot v_y.
 \label{eq:xor}
\end{equation}
The Gram formulation is a semidefinite program, with a quantum realization
available from the vector representation \cite{cleve2004}. This is an exact
characterization of the unrestricted two-party XOR value; a numerical solve
remains approximate. It does not extend unchanged to arbitrary multipartite
XOR tests or arbitrary non-XOR predicates.

\paragraph{Quantum witnesses.}
An explicit density matrix and local POVMs define a feasible strategy by
Eq.~\eqref{eq:quantum}. Check dimensions, Hermiticity, positivity, unit trace,
and POVM completeness before evaluating it. Floating-point checks yield a
numerically validated witness. Symbolic exact checks or interval error bounds
are needed for certified feasibility. Alternating optimization can search for
witnesses in a fixed dimension, but its stationary point is a lower-bound
candidate, not a global optimum.
The implemented bipartite see-saw alternates a largest-eigenvector state
update with local POVM semidefinite updates and uses seeded restarts.

\subsection{Symbolic certificates: CHSH as a complete example}
Use observables $A_0,A_1,B_0,B_1$ with $A_x^2=B_y^2=\id$ and
$[A_x,B_y]=0$. Let
$\mathcal B=A_0(B_0+B_1)+A_1(B_0-B_1)$. Direct expansion yields
\begin{equation}
 2\sqrt2\id-\mathcal B
 =\frac1{\sqrt2}\left[
 \left(A_0-\frac{B_0+B_1}{\sqrt2}\right)^2+
 \left(A_1-\frac{B_0-B_1}{\sqrt2}\right)^2\right].
 \label{eq:chsh-sos}
\end{equation}
Each squared Hermitian operator is positive, proving
$\langle\mathcal B\rangle\leq2\sqrt2$. The state
$\ket{\Phi^+}=(\ket{00}+\ket{11})/\sqrt2$, observables $A_0=Z,A_1=X$,
$B_0=(Z+X)/\sqrt2,B_1=(Z-X)/\sqrt2$, attain the bound. For uniform questions
and winning rule $a\oplus b=xy$,
\[
 \val_{\Lset}=\tfrac34,\qquad
 \val_{\Q_{\rm fin}}=\tfrac12+\tfrac{\sqrt2}{4},\qquad
 \val_{\NS}=1.
\]
The classical bound follows because the four required parity equations have
inconsistent XOR and three can be satisfied. The non-signalling value is
attained by $p(a,b\mid x,y)=1/2$ when $a\oplus b=xy$, and zero otherwise.
Equation~\eqref{eq:chsh-sos} proves the quantum upper bound independently of
solver tolerances. SOS methods also expose the operator relations needed for
self-testing, although extraction requires further arguments
\cite{bamps2015}.

\subsection{Metrics beyond one optimal value}
\label{sec:benchmarks}
The following exact literature values fix the conventions of the benchmark
constructors. They are mathematical expectations, not a claim that every
quantum value is certified by the implemented numerical backends.
\begin{center}
\begin{tabular}{lccc}
\toprule
Verifier and sampling & $\val_{\Lset}$ & $\val_{\Q_{\rm fin}}$ & $\val_{\NS}$\\
\midrule
Uniform CHSH & $3/4$ & $(2+\sqrt2)/4$ & $1$\\
Three-party GHZ, uniform even inputs & $3/4$ & $1$ & $1$\\
Magic square, uniform row/column & $8/9$ & $1$ & $1$\\
Three-party GYNI, uniform even inputs & $1/4$ & $1/4$ & $1/3$\\
\bottomrule
\end{tabular}
\end{center}
GHZ requires $\sum_i a_i=\frac12\sum_i x_i\pmod2$ under the even-input
promise; the three-party case of the parity-game analysis gives its values
\cite{brassard2003}. Magic square uses an even row for Alice, an odd column
for Bob, and agreement at their intersection \cite[Sec.~5]{brassard2005}.
This is the row/column version, not the row-or-column/single-cell variant
whose classical value is $17/18$ \cite[Sec.~3.3]{cleve2004}.
The perfect quantum strategies imply the stated non-signalling values.
GYNI accepts iff $(a_1,a_2,a_3)=(x_2,x_3,x_1)$. Almeida et al. prove that
quantum resources give no advantage for any input prior, and derive the
non-signalling value $1/3$ for this specific even-input prior
\cite[Eqs.~(10)--(12)]{almeida2010}. GYNI therefore tests an important failure
mode: a non-signalling gap alone does not demonstrate a quantum advantage.

For normalized games, report the additive gap
$\Delta=\val_Q-\val_L$ and the actual score conventions. Ratios are unstable
under harmless offsets and should not be compared without a normalization.
For a specified noise model
$p_v=vp_*+(1-v)p_0$ with $S_G(p_*)>\val_L>S_G(p_0)$, the visibility threshold
for violating the local bound is
\[
 v_{\rm crit}=\frac{\val_L-S_G(p_0)}{S_G(p_*)-S_G(p_0)}.
\]
This is a threshold for that family and witness, not for all noisy states or
all possible measurements. Loss, biased inputs, misalignment, and adversarial
noise require separate models. Conditioning on successful detection can change
the local bound and cannot be treated as unconditional performance.

Other useful questions include dimension/entanglement witnesses, nonlocal
membership, distance to the local set, robust extraction, randomness against
quantum side information, and sample complexity. These require different
optimization objects and assumptions. A game-value gap alone supplies neither
a device-independent entropy bound nor cryptographic security. Likewise,
finite-sample significance requires a sampling and memory model, explicit
confidence levels, and a valid statistical test; point estimates are not
confidence bounds.

\subsection{Why no complete automatic analyzer can exist}
There are finite games whose entangled value is not attained in finite
dimension \cite{slofstra2019}. More decisively, the gap problem distinguishing
entangled value $1$ from value at most $1/2$ is undecidable in general
\cite{ji2020}. A universal finite game language is therefore compatible with
the impossibility of a uniformly terminating, complete value-analysis
algorithm. The engineering consequence is typed partial procedures: exact
algorithms on restricted fragments, certified bounds when available, numerical
evidence with explicit limits, and an honest unresolved interval otherwise.

\section{Self-testing is a typed extraction claim}
\label{sec:selftest}

Classical input/output data can constrain untrusted quantum apparatus up to
unavoidable freedoms. The original self-testing formulation compares devices
with a reference realization through local changes of representation
\cite{mayers2004}; modern definitions distinguish state and measurement
certification and quantify robustness \cite{supic2020}.

\subsection{State extraction and the right quantifiers}
Fix a resource class of tensor-product strategies, a game $G$, a target pure
state $\ket\phi$, and a threshold $u$. Define target overlap after local
extraction by
\begin{equation}
 F_{\rm ext}(\rho,\phi)=
 \sup_{\Lambda_1,\ldots,\Lambda_n}
 \bra\phi(\Lambda_1\otimes\cdots\otimes\Lambda_n)(\rho)\ket\phi,
 \label{eq:extraction}
\end{equation}
where each $\Lambda_i$ is a local completely positive trace-preserving map
into the corresponding target system. A robust state-extraction statement is
\begin{equation}
 \forall (\rho,M)\in\mathcal R:\quad
 S_G(p_{\rho,M})\geq u-\varepsilon
 \ \Longrightarrow\ F_{\rm ext}(\rho,\phi)\geq1-f(\varepsilon).
 \label{eq:robust}
\end{equation}
The channels can depend on the unknown realization; the target and function
$f$ cannot. An exact self-test has $u=\val_{\mathcal R}(G)$, $f(0)=0$,
and a nonvacuous optimal realization. A robust self-test requires
$f(\varepsilon)\to0$ as $\varepsilon\to0$. The fidelity convention in
Eq.~\eqref{eq:extraction} is the overlap, not its square root.

Local isometries give a representation-level description:
the physical state maps to a target tensor an unobserved auxiliary state.
Auxiliary entanglement, unused Hilbert-space sectors, and local basis choices
are permitted. Measurement certification additionally compares the action of
physical operators on the tested state with target operators under the same
local maps. Operator-norm equality on the entire physical Hilbert space is
generally unjustified because the experiment does not probe unused sectors.

Pairwise convertibility of two optimal physical realizations is therefore
not the definition of a unique self-tested target. Both may contain a common
extractable target with different auxiliary systems. To disprove extraction
of a proposed $\phi$, one must exhibit an admissible optimum for which
$F_{\rm ext}<1$, or establish another obstruction to the stated equivalence.
Distinct numerical state vectors, Schmidt spectra before removal of auxiliary
factors, or inequivalent measurement choices alone do not settle this question.

\begin{proposition}[A common target with incomparable realizations]
Let $\ket{\Phi_d}=d^{-1/2}\sum_{j=0}^{d-1}\ket{jj}$. The states
$\ket{\Phi_6}$ and $\ket{\Phi_{10}}$ both attain the optimal CHSH score, but
neither can be transformed into the other by a product of local trace-preserving
quantum channels. Both exactly extract $\ket{\Phi_2}$.
\end{proposition}
\begin{proof}
Relabeling a local basis factors $\Phi_6$ as $\Phi_2\otimes\Phi_3$ and
$\Phi_{10}$ as $\Phi_2\otimes\Phi_5$. Use the CHSH measurements on the first
factor and discard the second for extraction. If product local channels
transform a pure $\Phi_d$ into a pure $\Phi_m$, their local Stinespring
isometries transform it into $\Phi_m\otimes\eta$ for a pure environment state
$\eta$. Local isometries preserve Schmidt rank, hence $d=mr$ with
$r$ the Schmidt rank of $\eta$. Neither $6/10$ nor $10/6$ is an integer.
\end{proof}
This statement concerns product local channels without communication;
LOCC convertibility is a different relation.

\subsection{Complex conjugation and information loss}
For a standard classical-interface Bell experiment, simultaneous complex
conjugation of $\rho$ and all local POVM matrices preserves every probability:
the Born probability is real and conjugation conjugates its trace. Complex
conjugation is not in general a physical channel. Measurement targets with
complex entries must therefore specify the tolerated conjugation freedom;
flagged simulations can require more than a single unflagged choice
\cite{mckague2010}. A compiler must not promise unique complex apparatus
merely because a numerical optimizer chose one representative.

There are three distinct uniqueness questions. A scalar Bell score can expose
a unique behavior; a behavior can still admit many physical realizations;
and these realizations may or may not extract a common target under the
chosen equivalence. Strict convexity or uniqueness of an SDP moment matrix
does not establish the last property. A finite moment relaxation may also
contain points without a corresponding physical strategy.

\subsection{What can be automated responsibly}
An SOS certificate $u\id-\mathcal B=\sum_kP_k^\dagger P_k$ implies
$P_k\ket\psi=0$ at exact saturation. At score deficit $\varepsilon$,
\[
 \sum_k\|P_k\ket\psi\|^2\leq\varepsilon.
\]
These relations are useful inputs to a proof of an extraction map; they are
not themselves such a map. Analytic operator inequalities can bound extraction
fidelity directly for important games \cite{kaniewski2016}. A further route
is to choose an extraction ansatz and lower-bound its fidelity through moment
relaxations, but the ansatz, operator identities, and bound direction must all
be checked.

The initial toolchain verifies a CHSH SOS identity and evaluates reference
strategies. It does not claim to derive or verify a general self-testing
theorem. A future certificate format should record the target, allowed
dimensions, extraction maps or a proved construction, gauge conventions,
distance metric, score normalization, robustness function, and a proof
dependency list. This makes ``self-testing supported'' a precise capability
rather than a boolean attached to a game name.

\section{Implementation architecture and abstraction cost}
\label{sec:architecture}

\subsection{A small core with specialized analysis paths}
The first executable fragment consists of finite one-round games with
classical questions/answers and fixed per-party alphabet sizes. Rational
sampling and payoff tables have exact semantics. A resource specification
belongs to an analysis request, so the same verifier can be compared under
several resource models. Unsupported model/backend combinations should fail
explicitly rather than silently weaken the assumptions.

The implementation separates four layers: constructors and semantic
validation; a finite game IR and serialization; algebraic transformations;
and analysis backends with evidence metadata. Examples exercise parity and
consistency rules rather than serving as separate solver implementations.
The finite table is deliberately small enough to inspect and test. Structured
and lazy IR nodes are an extension point, not a completed optimizing compiler.

Python is the initial host language because exact rationals, sparse linear
algebra, symbolic manipulation, and convex modeling can interoperate without
cross-language serialization during exploratory work. NumPy/SciPy handle
numeric arrays and linear programs; SymPy supplies exact algebra; CVXPY
provides solver-independent semidefinite modeling. Solver selection is a run
configuration, not a mathematical resource label. Open-source solvers provide
a portable baseline and commercial solvers can be optional.

The analogy with tile-oriented GPU DSLs is useful at the design level: retain
structure, validate contracts, and lower to specialized kernels only when
needed. Nonlocal-game analysis has different dominant costs. Exponential
strategy spaces and semidefinite matrix dimensions cannot be eliminated by
removing language overhead. A GPU compiler, a new parser, or a Rust/C++ core
should be introduced only after profiling identifies a repeated bottleneck
and the mathematical IR is stable. Existing packages such as Inflation offer
specialized causal-analysis functionality that can later be integrated
through adapters \cite{boghiu2023}.

\subsection{Implemented fragment versus extension roadmap}
The initial analysis fragment covers exact deterministic local optimization,
non-signalling linear optimization, two-party XOR semidefinite optimization,
optional rational non-signalling primal/dual certification,
selected bipartite moment relaxations, explicit finite-dimensional quantum
strategy evaluation, and a symbolic CHSH certificate. Bipartite
fixed-dimension see-saw searches are also available. Mixture, relabeling,
and parallel-product constructors operate on the finite IR. Numerical results
remain labeled as numerical; the exact local result and symbolic certificate
have separate evidence mechanisms.

Section~\ref{sec:resourceextensions} adds composable state and measurement
restrictions, fixed-state and PPT-constrained searches, independent-source
witness evaluation, and a nonlinear bilocality test. Section~\ref{sec:networkcompatibility}
implements exact source-disjoint marginal checks, triangle event inequalities,
and constructive membership for the no-input bilocal chain. General network
compatibility and global quantum network optimization remain outside this scope.

Full NPA level generation, symmetry reduction, sparse/lazy
compiler passes, general network optimization and inflation adapters, trusted quantum interfaces,
sequential combs, general extraction certificates, and finite-sample
cryptographic analysis are future work. In particular, documenting a network
resource does not make an ordinary Bell NPA backend a network optimizer.
Implementation/API documentation is the authority for the exact callable
surface; this report specifies its mathematical contract.

\subsection{Costs and validation obligations}
If the finite table has $N=|X||A|$ entries, constructing it is at least linear
in $N$. Exact local search involves exponentially many deterministic response
functions even with best-response elimination. A moment relaxation with $m$
selected words stores a matrix with $m^2$ entries; generic dense cone
operations can dominate runtime. Product constructors can multiply arena
sizes before any solver starts. The frontend should estimate these sizes and
permit a caller-specified budget before materialization.

An efficient lowering can exploit a parity structure through
Eq.~\eqref{eq:xor}, sparse payoff support, symmetry with a justified averaging
argument, or repeated symbolic subexpressions. A symmetry average requires
the feasible model to be convex and invariant; applying it indiscriminately
to an independent-source model can invalidate source independence.

Validation has two levels. Mathematical regression tests compare known exact
benchmarks, transformation identities, and witness scores. Implementation
tests check malformed distributions, shape mismatches, unsupported resources,
serialization, and numerical residual reporting. Random tiny games can
cross-check the optimized classical backend against direct enumeration, and
independent solvers can expose numerical sensitivity. Such agreement is
valuable evidence, but does not promote approximate SDP values to exact
certificates.

Reproducibility means declared dependencies, portable commands, public source
citations, and generated artifacts kept outside versioned source. Run reports
may record solver/package versions, tolerances, random seeds, and timings;
the library should not encode a particular workstation, solver license, or
environment path. Downloaded articles can be cached locally, while citations
must use stable public identifiers.

\section{Research program and falsifiable next steps}
\label{sec:research}

The framework's contribution at this stage is a precise separation of
semantics, compilation contracts, and evidence types, with elementary
preservation results. It does not yet establish a new separation theorem or
a complete basis of all nonlocal games. The following questions would turn
the architecture into a substantive research program.

\paragraph{A useful basis of constructors.}
Express benchmark families using parity, consistency, sampling, local wiring,
hidden/public selection, and parallel composition. Measure description size,
lowered size, compilation cost, and analysis cost separately. Investigate
which additional constructors are needed for linear-system, graph-coloring,
and synchronous games. Universality through arbitrary lookup tables is easy;
a basis that preserves exploitable structure is the meaningful target.

\paragraph{Certificates that compose.}
Start with component value bounds and SOS relations. Use the deficit rule of
Section~\ref{sec:composition} to track robustness under hidden mixtures, then
identify sufficient compatibility conditions for the component extraction
maps. A theorem should quantify the cost in score deficit, local dimension,
and target distance. Merely adding two self-testing scores is not a proof
that their certified subsystems coexist as a tensor product.

\paragraph{Resource-aware optimization.}
Construct pairs of arenas with identical payoff syntax but distinct source,
dimension, or trust assumptions. Ensure that the checker rejects invalid
backend reuse and that every accepted relaxation has the documented bound
direction. The no-input obstruction in Proposition~\ref{prop:noinput} is a
particularly useful test: a network compatibility problem should not be
mistaken for an ordinary linear game-value separation.

\paragraph{Certified numerical discovery.}
Search over structured game parameters, producing a candidate strategy lower
bound and a local upper bound for separation. Then seek a quantum upper bound
and reconstruct an exact dual or SOS certificate. Keep candidate discovery
and proof checking separate. Optimizing gaps over unconstrained score tables
is ill-posed because scaling increases the gap; normalize sampling/payoff and
control alphabet size, coefficient complexity, and physical noise assumptions.

\paragraph{A staged benchmark suite.}
CHSH tests value conventions, XOR specialization, an explicit witness, and
an exact SOS identity. GHZ parity tests expose multipartite semantics.
Magic-square consistency tests exercise nonbinary answer alphabets and
perfect quantum strategies. Degenerate zero-weight questions, biased
sampling, negative general payoffs, and no-input games test boundary cases.
Subsequent stages should add a dimension witness, a nonlinear network
compatibility problem, a semiquantum game, and a genuinely sequential
protocol. Each addition should come with an explicit semantic extension and
a demonstrably valid analysis method.

\paragraph{Open design question.}
Can a modest typed causal IR support all of these fragments while keeping
resource-preservation proofs small and backend-specific? This is a research
question about a useful representation, not a conjecture that every game
admits efficient solution. Success can be measured by reuse, verified
lowerings, smaller descriptions, and reproducible certificates even when
many values remain unresolved.

\section{Composable restrictions and independent-source witnesses}
\label{sec:resourceextensions}

A resource description becomes more expressive when restrictions can be
combined without erasing their individual meanings. The finite-interface
extension records a base correlation model, local dimension bounds, an
external shared-randomness policy, a state constraint, a measurement class,
and optionally a fixed supplied state. An admissible realization must satisfy
all applicable restrictions. These fields do not imply that a backend exists
for every combination.

\subsection{Dimensions, randomness, and admissible operations}
A local dimension bound $d_i$ limits the Hilbert space on which player $i$'s
state and measurement act. A fixed state supplies its own tensor-factor
ordering and dimensions; a larger numerical bound does not automatically
append an ancilla or replace the fixed state. In particular, a restriction to
projective measurements means projectors on the declared system. General
POVMs cannot be silently replaced by projective dilations when dimension or
ancilla access is constrained. Dimension-sensitive outer relaxations require
additional machinery beyond ordinary NPA \cite{navascues2015}.

External shared randomness allows the players to choose a common admissible
strategy branch before receiving their questions. With a fixed-state
contract, every branch uses the same supplied state; only otherwise permitted
choices may vary. This policy is separate from the correlations already
present in a mixed quantum state. Forbidding external randomness does not
require that the density matrix be pure, nor does it remove classical
correlations encoded inside the specified quantum resource.

\begin{proposition}[External convexification and linear scores]
\label{prop:convexification}
For any nonempty admissible behavior set $\mathcal R$ and linear score $S_G$,
\[
 \sup_{p\in\operatorname{conv}\mathcal R}S_G(p)
 =\sup_{p\in\mathcal R}S_G(p).
\]
Consequently external shared randomness changes neither the optimal linear
score nor a valid linear upper bound, although it can change membership and
nonlinear witnesses.
\end{proposition}
\begin{proof}
The score of a finite convex mixture is the same mixture of the branch
scores and cannot exceed their supremum. The reverse inequality follows from
$\mathcal R\subseteq\operatorname{conv}\mathcal R$.
\end{proof}
This observation does not authorize removing the randomness policy from a
resource type. A membership query, fixed-dimensional realization, or network
compatibility test can distinguish a set from its convex hull.

\begin{proposition}[Product resources have the local linear value]
\label{prop:productvalue}
Suppose the quantum state is a tensor product across all players, either fixed
or freely chosen. Allow arbitrary local POVMs, or arbitrary projective
measurements with zero projectors permitted. Then every finite game has the
same optimal score as its ordinary classical local value, with or without
external shared randomness.
\end{proposition}
\begin{proof}
The Born probabilities factor across players, so every allowed behavior is
local. Conversely, any deterministic response is implemented by assigning
$\id$ to its selected outcome and zero to every other outcome, independently
for each local question. These effects form a projective measurement on any
nonzero local system. Such responses attain the local linear optimum.
\end{proof}
The statement concerns values, not equality of unconvexified behavior sets.
It also depends on permitting degenerate measurements: a requirement that all
outcomes have nonzero rank would define a different resource model.
The implementation uses this exact reduction only for a symbolic product-state
resource without a supplied numerical state. A fixed matrix accepted as
approximately product within a tolerance is not thereby proved exactly product;
its numerical evaluation or search must retain a numerical evidence label.

\subsection{Fixed states and positive partial transpose}
For a supplied state $\rho_0$, the optimization variables are only the allowed
local measurements. Alternating measurement optimization holds $\rho_0$
unchanged. A state eigenvector update would solve a different problem. Even
when each measurement update is an SDP, the combined optimization is
nonconvex and a see-saw output supplies only numerical achievable evidence.

For a nonempty proper subset $S$ of parties, let $T_S$ denote partial
transposition on those tensor factors. A PPT restriction across a specified
collection $\mathcal C$ of cuts is
\[
 \rho\succeq0,\qquad \Tr\rho=1,\qquad
 \rho^{T_S}\succeq0\quad(S\in\mathcal C).
\]
Complementary cuts impose equivalent positivity conditions. The chosen cuts
must nevertheless be explicit; being PPT across one cut is not a declaration
of positivity across every bipartition. Every state separable across a cut
is PPT across that cut \cite{peres1996}. The converse holds for bipartite
$2\times2$ and $2\times3$ systems, but fails in general
\cite{horodecki1996}. Thus PPT is a tractable constraint, not an alias for
separability or for absence of Bell nonlocality.

\begin{proposition}[Scope of a fixed-measurement PPT optimization]
Fix local measurements $M$ and their Bell operator $\mathcal B_M$. The program
\[
 u_{\rm PPT}(M)=\max_\rho\Tr(\rho\mathcal B_M)
 \quad\text{subject to the stated PPT constraints}
\]
is an SDP. Its exact optimum bounds the score of every state obeying those
constraints with these same measurements. It is not in general an upper bound for
$\sup_{M'}u_{\rm PPT}(M')$.
\end{proposition}
\begin{proof}
The objective is linear in $\rho$, and positivity, trace, and partial
transpose positivity are semidefinite or affine constraints. The fixed
measurement family is one choice in the larger optimization, so
$u_{\rm PPT}(M)\leq\sup_{M'}u_{\rm PPT}(M')$.
\end{proof}
A see-saw that alternates this state problem with local POVM updates remains
a heuristic for the joint problem. Numerical repair must respect all active
cones, and scores must be recomputed after repair. For example, mixing a
Hermitian normalized candidate with the maximally mixed state can move both
its ordinary and partial-transpose spectra into the positive cone. Finite
precision checks and repairs are not exact certificates.

\subsection{Networks retain the source factorization}
An independent-source witness specifies a hypergraph and one state or latent
variable distribution for each source. In a quantum witness, the global
source state is constructed as $\bigotimes_e\rho_e$. Its factors are then
permuted into party order before local joint POVMs are applied. Source order,
endpoint order, and the resulting party-local factor order are part of the
representation. Accepting an arbitrary global density matrix in place of
these factors would discard independence.

A classical finite-source witness similarly uses an explicit product of
rational latent-variable distributions and rational local response tables.
Such a witness can be evaluated exactly. It is a constructive compatibility
certificate for its resulting behavior, not a complete membership procedure
for arbitrary supplied behavior tables. Quantum source witnesses are checked
numerically. General network outer approximations, for example inflation,
remain a separate backend capability \cite{wolfe2021}.

\begin{proposition}[Classical network convex hull and linear values]
\label{prop:networklinear}
Fix any independent-source classical network, with unrestricted local response
functions. Let $\mathcal N$ be its behavior set in a finite Bell arena and
$\Lset$ the ordinary local polytope on the same arena. Then
\[
 \operatorname{conv}\mathcal N=\Lset,\qquad
 \sup_{p\in\mathcal N}S_G(p)=\val_{\Lset}(G).
\]
The result also holds if source distributions and finite latent cardinalities
are fixed, provided they are valid and every party can ignore its sources.
\end{proposition}
\begin{proof}
A network realization is ordinary local by collecting all source variables
into one hidden variable with their joint product distribution. Conversely,
every deterministic local response tuple belongs to $\mathcal N$: the
parties ignore their sources. The convex hull of these deterministic tuples
is $\Lset$. The value equality follows either from these inclusions or from
Proposition~\ref{prop:convexification}.
\end{proof}
Thus an exact ordinary classical optimizer is a sound network \emph{linear
value} backend under these assumptions, with an optimal witness that ignores
all sources. It is not a sound network \emph{membership} backend. Source
constraints can distinguish behaviors and nonlinear witnesses even though
no linear functional separates their optimum values. Restricting local
response families or imposing additional observed statistics changes the
premises and requires a new analysis.

The bilocal chain $A$--$B$--$C$ has independent sources on $AB$ and $BC$.
Let $A$ and $C$ have binary inputs $x,z$ and binary outputs, while $B$ has
one input and a two-bit output $(b_0,b_1)$. With signs
$A_x=(-1)^a$, $B^k=(-1)^{b_k}$, and $C_z=(-1)^c$, define
\[
 I=\tfrac14\sum_{x,z}\langle A_xB^0C_z\rangle,\qquad
 J=\tfrac14\sum_{x,z}(-1)^{x+z}\langle A_xB^1C_z\rangle.
\]
Every classical bilocal behavior obeys
\begin{equation}
 \sqrt{|I|}+\sqrt{|J|}\leq1.
 \label{eq:bilocal}
\end{equation}
This is a nonlinear distribution witness \cite{branciard2012}.
For completeness, write $a_\pm(\lambda)=(A_0(\lambda)\pm
A_1(\lambda))/2$ and similarly $c_\pm(\nu)$. Independence gives
$|I|\leq uv$ and $|J|\leq u'v'$, where
$u=\mathbb E|a_+|$, $u'=\mathbb E|a_-|$, and similarly $v,v'$.
Because $|a_+|+|a_-|\leq1$, Cauchy--Schwarz yields
$\sqrt{uv}+\sqrt{u'v'}\leq\sqrt{(u+u')(v+v')}\leq1$.

Two independent Bell pairs, Bob's joint measurement of $Z\otimes Z$ and
$X\otimes X$, and the end measurements
$(Z+(-1)^xX)/\sqrt2$ give $I=J=1/2$. The resulting value is $\sqrt2$.
This construction tests source factorization and tensor ordering as well as
the nonlinear witness. It is not an ordinary scalar game-value claim.

The same expression exhibits why global shared randomness cannot be added
silently. One deterministic behavior has all output signs $+1$, giving
$(I,J)=(1,0)$. Another uses $A_x=(-1)^x$, $C_z=(-1)^z$, and Bob's signs $+1$,
giving $(I,J)=(0,1)$. Both are bilocal. Their equal mixture gives
$(I,J)=(1/2,1/2)$ and violates Eq.~\eqref{eq:bilocal}. The mixture is
ordinary Bell-local, but needs a common branch variable unavailable in the
independent-source model. The ordinary local polytope is therefore an invalid
replacement for bilocal compatibility.

\subsection{Capabilities and evidence remain separate}
Resource validation asks whether a supplied strategy satisfies a contract.
Search asks for a high-scoring witness inside that contract. Relaxation asks
for an outer bound, and certification independently checks a mathematical
claim. These operations must not share an ambiguous success flag. In
particular, a source witness can be flattened for ordinary Born-rule
evaluation while its source structure is retained for checking; that
flattening does not justify using an unrestricted optimizer as a network
optimizer.

Dimension-sensitive global bounds, generic separability or entanglement-depth
constraints, network optimization and inflation adapters, trusted quantum
questions, and sequential instruments require further backends. The resource
schema makes these extension points explicit without asserting that a single
solver can handle them all.

\section{Nonlinear compatibility checks with explicit scope}
\label{sec:networkcompatibility}

The value equality in Proposition~\ref{prop:networklinear} makes a
behavior-valued task essential for testing source independence. A
compatibility checker must report either a constructive realization, a
violated necessary condition, or an unresolved case. Satisfying a finite
list of necessary conditions does not establish compatibility.

\subsection{Independence of source-disjoint marginals}
For a set $U$ of parties, let $E(U)$ denote the union of sources incident to
at least one party in $U$.
\begin{proposition}[Source-disjoint groups]
\label{prop:sourcedisjoint}
Suppose disjoint party groups $U,V$ satisfy $E(U)\cap E(V)=\varnothing$.
In an independent-source classical or tensor-product quantum network,
with no additional global shared randomness or communication,
\[
 p(a_U,a_V\mid x)=p(a_U\mid x_U)\,p(a_V\mid x_V).
\]
The equation is a marginal over all other outputs, not a condition on their
values. Inputs are fixed independently of the sources.
\end{proposition}
\begin{proof}
In the classical model, summing other local outputs removes their normalized
response factors. The remaining response factors depend on disjoint
collections of independent source variables, so their expectation factors.
In the quantum model, summing the other POVMs gives identities. Tracing their
systems out of the product source state leaves a product between the systems
in $U$ and those in $V$, since no source meets both groups. The Born trace
therefore factors. Independence from inputs outside each group follows in
the same way.
\end{proof}
In a bilocal chain this forces $p(a,c\mid x,z)=p(a\mid x)p(c\mid z)$.
Conditioning on Bob's output can correlate the end parties through
entanglement swapping, so the unconditional factorization must not be
applied to such postselected data. An observed violation excludes both the
classical and the specified quantum independent-source model, under the
stated assumptions.

\subsection{A triangle event inequality}
For the triangle with independent pair sources $AB$, $BC$, and $CA$, every
classical behavior obeys, at each fixed input tuple,
\begin{equation}
 p(a,b,c\mid x,y,z)^2\leq
 p_A(a\mid x)\,p_B(b\mid y)\,p_C(c\mid z).
 \label{eq:trianglefinner}
\end{equation}
This is the triangle specialization of Finner's product-space inequality
\cite{finner1992}. Network versions and a quantum extension for independent
bipartite sources are established by Renou et al.~\cite{renou2019limits}.
The checker records applicability to both classical and quantum independent
pair sources. The elementary proof below establishes the classical case;
the quantum case uses the cited theorem rather than replacing operator
algebra by classical hidden variables.

Let $r,s,t$ be the independent variables of the three sources and let
$f(r,t)$, $g(r,s)$, and $h(s,t)$ be the local conditional probabilities of
the selected events. Private local randomness is already absorbed into
these functions, so $0\leq f,g,h\leq1$. At the fixed inputs,
$p(a,b,c)=\mathbb E[fgh]$. Two Cauchy--Schwarz applications give
\begin{align*}
 \mathbb E[fgh]
 &\leq \mathbb E_{s,t}\!\left[h(s,t)
       \sqrt{\mathbb E_r f(r,t)^2\,\mathbb E_r g(r,s)^2}\right],\\
 \bigl(\mathbb E[fgh]\bigr)^2
 &\leq \mathbb E[h^2]\,
   \mathbb E_{s,t}\!\left[\mathbb E_r f(r,t)^2\,
                         \mathbb E_r g(r,s)^2\right]\\
 &=\mathbb E[f^2]\,\mathbb E[g^2]\,\mathbb E[h^2]
 \leq\mathbb E[f]\,\mathbb E[g]\,\mathbb E[h].
\end{align*}
The last expectations are the three marginal event probabilities.

For example, the common-random-bit distribution
$p(000)=p(111)=1/2$ violates Eq.~\eqref{eq:trianglefinner}:
the event $000$ has squared probability $1/4$, while the product of its
marginals is $1/8$. Thus a common tripartite classical source cannot be
silently replaced by three independent pair sources. A nonviolation of
these event inequalities supplies no general realization theorem. The
checker does not interpret an arbitrary nonsignalling table as a network
of independent nonsignalling resources; that is a separate physical model.

\subsection{An exactly solvable no-input chain fragment}
When all parties in a bilocal chain have one input, source-disjoint
independence is also sufficient for compatibility, for arbitrary finite
output alphabets \cite{renou2019limits}.
\begin{proposition}[No-input bilocal reconstruction]
\label{prop:noinputbilocal}
A normalized distribution $p(a,b,c)$ admits a classical independent-source
realization on $A$--$B$--$C$ if and only if
$p_{AC}(a,c)=p_A(a)p_C(c)$. The same condition characterizes the
finite-dimensional quantum independent-source behaviors in this no-input
fragment.
\end{proposition}
\begin{proof}
Necessity in both models follows from
Proposition~\ref{prop:sourcedisjoint}. For sufficiency, the first source
samples $\lambda=a$ from $p_A$, and the second samples $\nu=c$ from $p_C$,
independently. Alice and Charlie output these variables. Bob receives both
and samples
\[
 q(b\mid\lambda=a,\nu=c)=
 \frac{p(a,b,c)}{p_A(a)p_C(c)}
\]
whenever the denominator is positive. The assumed factorization normalizes
this response. If the denominator is zero, normalization and nonnegativity
force every $p(a,b,c)$ in that row to vanish, and Bob may use any normalized
response there. The constructed joint probability is exactly $p(a,b,c)$.
A finite classical source/response model embeds into quantum theory using
orthogonal classical registers and diagonal POVMs.
\end{proof}
Rational input data produce rational source probabilities and responses, so
both the independence condition and the reconstructed witness can be checked
exactly. The proposition does not extend to multi-input bilocal games: one
must then use a common source model across the input choices, and the
marginal equality alone is insufficient.

\subsection{Certificates, tolerances, and unresolved cases}
For rational tables, a violation certificate can name an input profile, an
output event or marginal event, and the exact rational residual. A checker
recomputes that residual from the full supplied table and source graph,
including normalization and all applicability conditions. No floating square
root is needed for Eq.~\eqref{eq:trianglefinner}. Such a certificate excludes
a resource model for the supplied distribution; it is not an experimental
confidence statement.

The general compatibility checks accept only exact rational, normalized,
non-signalling tables and reject floating input. The separate bilocal
inequality evaluator also supports floating tables with explicit tolerances
and numerical evidence labels. No violation found means only that the
implemented conditions did not exclude the data. The constructive no-input chain result is an exception with a
proved sufficient condition and an explicit witness, not a precedent for
promoting generic passing checks to membership certificates.

\section{Multipartite strategy search and binary projective updates}
\label{sec:multipartitesearch}

For a finite-dimensional one-round tensor-product realization, alternating
optimization can update one local measurement at a time at any number of
parties. It does not need a bipartite restriction. With all other measurements
and the state fixed, the objective for party $i$ at question $x_i$ is
\[
 \sum_{a_i}\Tr(E^{(i)}_{a_i\mid x_i}H^{(i)}_{a_i\mid x_i}),
 \qquad E^{(i)}_{a_i\mid x_i}\succeq0,\quad
 \sum_{a_i}E^{(i)}_{a_i\mid x_i}=\id.
\]
The effective Hermitian matrices $H$ contract the state with the other
parties' effects and sum over their questions and answers with the game
coefficients. Their construction must retain the tensor-factor ordering,
including complex conjugation conventions implicit in the trace.

To make the convention explicit, permute party $i$ first and group the
remaining systems into one factor. Writing the state as
$\rho_{u r,v s}$ and the product of the other parties' effects as $O_{s r}$,
the contribution to the effective matrix is
\[
 H_{u v}=\sum_{r,s}\rho_{u r,v s}O_{s r}.
\]
Then the score contribution is $\Tr(EH)$. This contraction also covers one
party by using a one-dimensional empty tensor product. Tests compare the
optimized block against a separately assembled full-tensor SDP for complex
states and unequal local dimensions.

\begin{proposition}[Binary measurement block]
For Hermitian $H_0,H_1$, the maximum of
$\Tr(E_0H_0)+\Tr(E_1H_1)$ over $E_0,E_1\succeq0$ with
$E_0+E_1=\id$ is attained by a projective measurement:
\[
 E_0=\Pi_{>0}(H_0-H_1),\qquad E_1=\id-E_0.
\]
A choice of subspace in the zero eigenspace is immaterial to the value.
\end{proposition}
\begin{proof}
Eliminating $E_1$ leaves the constant $\Tr H_1$ plus
$\Tr(E_0\Delta)$ with $0\preceq E_0\preceq\id$ and
$\Delta=H_0-H_1$. In an eigenbasis of $\Delta$, each diagonal entry of
$E_0$ is in $[0,1]$. Set it to one on positive eigenvalues and zero on
negative ones. The indicated projector realizes that optimum.
\end{proof}
The identity includes zero and full-rank projectors. A one-outcome
measurement is necessarily the identity. Thus a dimension-constrained
binary projective update needs no Naimark dilation and no measurement SDP.
The general multi-outcome projective constraint is nonconvex and is not
silently replaced by a POVM optimization.

The complete search remains nonconvex: optimality of each block does not
imply global optimality of the joint state-and-measurement strategy. Fixed
states remain unchanged; free states use an eigenvector or a PPT-constrained
state block. Every resulting strategy is checked against its dimension,
state, and measurement contract. Numerical spectra and feasibility tolerances
retain numerical evidence labels, including in the analytic binary update.
Independent-source networks need factorized source-state updates and are not
covered merely by allowing more parties in this ordinary tensor-product search.

\section{Exact primal and dual certificates for XOR games}
\label{sec:xorcertificates}

For a bipartite binary-output parity score, let $e_{xy}$ and $o_{xy}$ be its
even- and odd-parity rewards. Define the rational constant and weight matrix
\[
 c_0=\sum_{xy}\mu(x,y)\frac{e_{xy}+o_{xy}}2,
 \qquad W_{xy}=\mu(x,y)\frac{e_{xy}-o_{xy}}2,
 \qquad C=\begin{pmatrix}0&W/2\\W^T/2&0\end{pmatrix}.
\]
The vector characterization discussed in Section~\ref{sec:analysis} has primal
objective $c_0+\langle C,G\rangle$, subject to $G\succeq0$ and
$G_{ii}=1$. The off-diagonal factor $1/2$ in $C$ avoids counting each cross
correlator twice. The usual SDP solve is numerical. Its equality dual can
instead seed an independently checked rational certificate.

\begin{proposition}[XOR dual certificate]
Let $z$ be a rational vector and suppose the rational symmetric matrix
$S=\operatorname{Diag}(z)-C$ is positive semidefinite. Then
\[
 \val_{\Q_{\rm fin}}(G)\leq c_0+\sum_i z_i.
\]
The same bound applies to commuting strategies for this bipartite XOR fragment.
\end{proposition}
\begin{proof}
For every feasible Gram matrix, positivity gives
$\langle S,G\rangle\geq0$. Since its diagonal is one,
$\langle C,G\rangle\leq\sum_i z_i$. The XOR vector characterization
includes the relevant correlations \cite{cleve2004}.
\end{proof}
This proof uses only feasibility of the supplied dual; an optimal numerical
status is neither required nor sufficient. Signed rewards, biased sampling,
and zero-probability questions are handled by the same rational formulas.

\subsection{An exact PSD checker}
A symmetric rational matrix can be checked by successive Schur complements.
A negative diagonal pivot disproves positivity. A zero pivot requires every
remaining entry in that pivot's row to be zero: otherwise the corresponding
$2\times2$ principal minor has negative determinant. A positive pivot $d$
reduces the problem to the remaining matrix minus $vv^T/d$. These steps are
necessary and sufficient, so exact rational arithmetic handles singular
matrices as well as positive definite ones. No floating eigenvalue or tolerance
enters certificate verification. Matrix-size budgets limit rational elimination
cost; they do not weaken the checked condition.

\subsection{Repairing numerical candidates}
Round a numerical dual into rationals and first test its slack exactly. If
necessary, increase every diagonal coordinate by the same rational amount.
The choice
\[
 \delta=\max\left(0,\max_i\left[\sum_{j\ne i}|C_{ij}|-z_i\right]\right)
\]
gives a symmetric diagonally dominant slack with nonnegative diagonal and
therefore a guaranteed feasible endpoint. A fixed number of rational bisection
steps, each using the exact PSD checker, can reduce this repair while retaining
a feasible upper endpoint. The resulting objective increases by
$(|X_A|+|X_B|)\delta$ for the retained shift; this cost is reported explicitly.

The resulting value is a certified upper bound, not necessarily the optimum.
A numerical Gram matrix does not supply a certified lower bound. One may
separately compare with an exact classical strategy or an independently verified
quantum realization. Serialized certificates bind to the canonical verifier
and store rational diagonals and the claimed bound. Rechecking recomputes the
matrix, positivity, and bound from the supplied game, independently of the
original solver and its tolerances.

\subsection{Attainment from rational Gram data}

\begin{proposition}
Let $G$ be a real rational positive semidefinite matrix with unit diagonal
in party/question order and rank $r$. The score $c_0+\langle C,G\rangle$ is
attained by a finite strategy of local dimensions $D=2^{\lfloor r/2\rfloor}$
on each side for the bipartite XOR objective.
\end{proposition}
\begin{proof}
Factor $G$ over the reals to obtain unit vectors $u_x,v_y$ in $\mathbb R^r$.
For $k=\lfloor r/2\rfloor$, take the $2k$ Pauli strings with a prefix of $Z$'s,
one $X$ or $Y$, and an identity suffix on $k$ qubits. For odd $r$ include
$Z^{\otimes k}$, using the scalar identity when $r=1$. These Hermitian
involutions $\Gamma_i$ anticommute pairwise and obey
$\Tr(\Gamma_i\Gamma_j)/D=\delta_{ij}$. For a real unit vector $u$ the
observable $A(u)=\sum_i u_i\Gamma_i$ squares to the identity. On a maximally
entangled state of dimension $D$ per party, $A(u_x)$ and $A(v_y)^T$ have
correlator $\Tr(A(u_x)A(v_y))/D=u_x\cdot v_y$. Substitution in the objective
gives the stated score. This is a finite Clifford realization of the XOR
vector correspondence \cite{cleve2004}.
\end{proof}

The checker verifies the Gram constraints, rank, and objective by rational
arithmetic without constructing the real factors or the $D$-dimensional
observables. The dimension is sufficient for this realization. It supplies
a lower bound for dimension-capped value claims when both caps are at least
$D$. A smaller cap leaves this certificate inconclusive.

For discovery, symmetrize a numerical candidate, fix the diagonal to one,
and round off-diagonal entries to rational grid points. If the resulting $R$
fails exact PSD, let $M=\max_i\sum_{j\ne i}|R_{ij}|$. Then $M>1$ and
$(1-t)R+tI$ is PSD at $t=1-1/M$ by diagonal dominance. Bounded bisection
with exact feasibility checks can reduce $t$. Compute the achieved score on
the final checked matrix. An input-dependent output flip changes the sign of
its node vector, so relabelling transports $G$ by a signed permutation
congruence. The same transformation on the diagonal dual variables only
requires the underlying node permutation.

\section{A certified classical spiral-inflation fragment}
\label{sec:spiralinflation}

The triangle Finner test does not exhaust the restrictions imposed by source
independence. A small classical inflation supplies a stronger, still tractable
compatibility test. We use a fixed six-variable spiral construction from the
inflation framework of Wolfe, Spekkens, and Fritz
\cite{wolfe2019classicalinflation}; it is not a complete inflation hierarchy.
Its resource contract is explicitly classical fanout.

\subsection{Copied sources and observable marginal constraints}
Let $U,V,W$ be independent classical sources on $AB$, $BC$, and $CA$.
For stochastic response functions, introduce independent private randomness
at each inflated response; it is suppressed in the following notation.
Take independent identically distributed copies $U_1,U_2,V_1,V_2,W_1,W_2$
and define
\[
\begin{array}{lll}
 A_1=f_A(U_1,W_1),&B_1=f_B(U_1,V_1),&C_1=f_C(V_1,W_1),\\
 A_2=f_A(U_2,W_1),&B_2=f_B(U_1,V_2),&C_2=f_C(V_1,W_2).
\end{array}
\]
Every classical triangle realization therefore induces a joint distribution
$q$ of these six outputs. The following five triple marginals of $q$ are
fixed by the original behavior $p$:
\begin{equation}
\begin{aligned}
 q_{A_1B_1C_1}(a,b,c)&=p_{ABC}(a,b,c),\\
 q_{A_2B_2C_1}(a,b,c)&=p_{AC}(a,c)p_B(b),\\
 q_{A_1B_2C_2}(a,b,c)&=p_{AB}(a,b)p_C(c),\\
 q_{A_2B_1C_2}(a,b,c)&=p_{BC}(b,c)p_A(a),\\
 q_{A_2B_2C_2}(a,b,c)&=p_A(a)p_B(b)p_C(c).
\end{aligned}
\label{eq:spiralmarginals}
\end{equation}
For example, $(A_2,C_1)$ has the original $AC$ marginal, and all its
ancestors are disjoint from those of $B_2$. This proves the second
factorization. The other factorizations follow from the same argument.
The five sets are the maximal ancestrally-independent-expressible sets of
this spiral. Their pair marginals need not be imposed separately.

Fix finite output alphabet sizes $(d_A,d_B,d_C)$ and let
$d=d_Ad_Bd_C$. Order joint assignments lexicographically as
$(A_1,B_1,C_1,A_2,B_2,C_2)$, and order marginal rows by the five blocks in
Eq.~\eqref{eq:spiralmarginals}, then lexicographically within each block.
The induced marginal program is
\begin{equation}
 Aq=b(p),\qquad q\geq0,
 \label{eq:spirallp}
\end{equation}
where $A$ is a $5d\times d^2$ integer matrix containing exactly $5d^2$
nonzero entries, all equal to one. Each column contributes one entry to each
marginal block. Normalization of the first block implies $\sum q=1$.
The implementation guards these three sizes before enumerating
assignments, and stores $A$ sparsely. For binary outputs there are 40 rows,
64 variables, and 320 nonzeros.

\begin{proposition}[Scope of the spiral marginal test]
Every no-input classical triangle behavior makes
Eq.~\eqref{eq:spirallp} feasible. Consequently infeasibility excludes that
resource model. Feasibility of this program alone does not certify triangle
compatibility.
\end{proposition}
\begin{proof}
The copied-source construction produces an admissible nonnegative normalized
$q$, and its marginals satisfy Eq.~\eqref{eq:spiralmarginals}. The converse
would require reconstructing the source model, which the marginal LP does
not enforce.
\end{proof}
The common joint distribution construction relies on classical variables
being reusable at multiple response nodes. It does not justify copying
quantum source systems or assuming joint outcomes for incompatible quantum
operations. This LP therefore provides no general quantum-network
incompatibility claim. Quantum inflation requires a different operator
construction.

\subsection{Exact infeasibility certificates}
A candidate dual vector $y$ is a valid contradiction precisely when
\begin{equation}
 A^{\mathsf T}y\geq0,\qquad y^{\mathsf T}b(p)<0.
 \label{eq:spiralfarkas}
\end{equation}
Indeed, a hypothetical nonnegative $q$ satisfying the marginal equations
would give $0\leq q^{\mathsf T}A^{\mathsf T}y=y^{\mathsf T}b(p)<0$.
Conversely, Farkas' lemma guarantees a separating vector for every infeasible
program. Bounding the coordinates of $y$ during numerical discovery, for
example by $|y_j|\leq1$, loses no strict separation because a Farkas vector
can be scaled. These discovery bounds are not part of certificate validity.

All input probabilities and reconstructed certificate coefficients are
rational. Verification rebuilds $A$ and $b(p)$ and checks every inequality
in Eq.~\eqref{eq:spiralfarkas} exactly, independently of the numerical
solver. A useful exact repair adds $\delta=-\min(A^{\mathsf T}y)>0$ to
every coefficient in the first marginal block whenever a candidate has
negative dual slack. Every slack then increases by $\delta$, and the dual
margin increases by $\delta\sum p_{ABC}=\delta$. The repaired vector is
retained only if its final exact margin is still negative.

A numerically reported infeasibility without a verified dual is not an
exact exclusion. A rational primal $q$ satisfying the equations certifies
only feasibility of the relaxation, so its compatibility status remains
\emph{not ruled out}. Numerical failures, missing or malformed witnesses, and unsuccessful
rational reconstruction retain an explicit unresolved status or numerical
evidence label. Numerical feasibility additionally requires checked residuals;
solver option values that would be silently reset are rejected.

\subsection{A W-type benchmark beyond the Finner checks}
For binary outputs, the event inclusion on the six inflated variables gives
\begin{equation}
 p_A(1)p_B(1)p_C(1)\leq p_{ABC}(000)
 +p_{AC}(11)p_B(1)+p_{AB}(11)p_C(1)+p_{BC}(11)p_A(1).
 \label{eq:spiralw}
\end{equation}
To see the inclusion, if all three outer variables are one, either all
inner variables are zero, or at least one of $A_1,B_1,C_1$ is one. The
latter possibilities lie in the corresponding expressible triple events
on the right. Taking probabilities and applying
Eq.~\eqref{eq:spiralmarginals} proves the inequality.

The distribution $p(100)=p(010)=p(001)=1/3$ has a zero right-hand side and
left-hand side $1/27$. This provides a simple integer dual certificate:
in the stated zero-based row order, set $y_0=y_{15}=y_{23}=y_{31}=1$,
$y_{39}=-1$, and every other coefficient to zero. Direct checking gives
$\min(A^{\mathsf T}y)=0$ and $y^{\mathsf T}b=-1/27$.
Meanwhile every nonzero event passes the triangle Finner test, since
$(1/3)^2\leq(1/3)(2/3)(2/3)$. The benchmark thus exercises
stronger classical compatibility information without changing its evidence
scope. It is an incompatibility target, not a demonstrated quantum advantage:
no quantum realization of this target is supplied by this example.

\subsection{A polynomial witness, not a linear Bell functional}
For fixed output alphabets, $A$ is independent of the tested distribution.
An exactly verified $A^{\mathsf T}y\geq0$ therefore defines the universal
classical-triangle inequality $\Phi_y(p)=y^{\mathsf T}b(p)\geq0$.
Choosing $y$ by optimizing against one distribution does not invalidate
this statement once structural dual feasibility has been independently
checked. However, $b(p)$ contains linear, quadratic, and cubic terms in the
original probabilities. The witness is linear in the lifted marginal
coordinates, but is generally a polynomial in $p$, not an ordinary Bell
score.

Exported numerical-discovery records bind the complete behavior,
source graph, output alphabets, and row-order schema through canonical
fingerprints. Checking the same coefficient vector on different data must
explicitly rebuild and evaluate $b(p)$ in the new context. If a future
relaxation introduces distribution-dependent coefficients in $A$ itself,
the preceding universal-polynomial interpretation requires a new proof;
a single-instance dual cannot then be reused without checking the changed
constraints.

\section{Claims, evidence admission, and discovery semantics}
\label{sec:claims}

A solver result and a mathematical conclusion are distinct objects. A result
may contain a numerical objective, a strategy candidate, or a dual certificate.
A claim additionally specifies the verifier, resource model, threshold or
target, and quantifiers. The verification verdict is \emph{proved},
\emph{disproved}, or \emph{unknown}. Unknown is an epistemic outcome of the
available checks, not a third truth value of the underlying proposition.
Execution outcomes such as unsupported assumptions, invalid evidence, and
budget exhaustion are recorded separately and never treated as disproofs.

The initial claim fragment concerns value suprema in the standard resource
chain of Eq.~\eqref{eq:chain}. It does not decide membership, perfect-strategy
attainment, or self-testing. In particular, a proof that the finite-dimensional
quantum supremum is one does not alone establish an attained perfect strategy
\cite{slofstra2019}. More restrictive dimension, fixed-state, network, PPT, and
measurement contracts remain explicit; an inference rule must justify any
transport involving them before the claim engine accepts it.

\begin{proposition}[Bound transport and threshold inference]
If $\mathcal R\subseteq\mathcal S$, a lower bound for
$\val_{\mathcal R}(G)$ is a lower bound for $\val_{\mathcal S}(G)$,
and an upper bound for $\val_{\mathcal S}(G)$ is an upper bound for
$\val_{\mathcal R}(G)$. Given checked intervals
$L_R\leq\val_{\mathcal R}(G)\leq U_R$ and
$L_S\leq\val_{\mathcal S}(G)\leq U_S$, the strict gap claim
$\val_{\mathcal S}(G)-\val_{\mathcal R}(G)>\varepsilon$ is proved if
$L_S-U_R>\varepsilon$, and disproved if $U_S-L_R\leq\varepsilon$.
\end{proposition}
\begin{proof}
Resource monotonicity gives
$\val_{\mathcal R}(G)\leq\val_{\mathcal S}(G)$. Subtracting the checked
interval endpoints gives
\[
 L_S-U_R\leq\val_{\mathcal S}(G)-\val_{\mathcal R}(G)\leq U_S-L_R.
\]
The stated sufficient conditions follow, including the distinction between
strict and non-strict threshold comparisons.
\end{proof}
An NS or NPA upper bound above a classical value therefore does not prove a
quantum advantage. It leaves that possibility open. A numerical candidate can
guide the next search, but cannot replace a certified lower witness in this
inference. Contradictory checked intervals are treated as a conflict requiring
an audit of the assumptions and checker, not resolved by preferring a solver.

\subsection{Artifact identity is not proof authority}
Evidence artifacts contain a versioned kind, an exact verifier identity, a
complete source-resource representation, and an immutable canonical data
payload. The verifier identity retains sampling and payoff separately and
ignores only the descriptive name. Such an identity is not a quotient under
local relabellings, affine score changes, or experimentally equivalent sampling.
Those transformations require checked maps and their own inference rules.

Content digests detect changes to an envelope, but anyone can construct an
envelope and recompute its digest. Consequently evidence admission independently
replays a registered check against the actual game. Deterministic response
tables establish only their achieved score; an upper classical bound requires
exhaustive replay or another proof. Rational NS primal/dual witnesses and XOR
duals establish exactly the bounds checked by their respective verifiers.
Numerical observations have no bound authority, even if a supplied result label
claims exactness. Registering a label in an enumeration checks vocabulary, not
mathematical truth.

The admission record retains checker identifiers, versions, source models,
bounds, and artifact references. A serialized assessment is a receipt rather
than a standalone proof: its claim definition and raw evidence must remain
available for rechecking. The trusted boundary includes the mathematical model
compiler, resource-inclusion rules, parsing, arithmetic, and checker code.
Independent rational verification does not by itself constitute proof-assistant
formalization, and a shared compiler error can affect discovery and checking.

\subsection{A bounded discovery workflow}
The first executable workflow admits only a small set of precise evidence
kinds and standard unrestricted resource models. Count and size budgets limit
artifact admission, exhaustive replay, and matrix construction; they do not
promise a global wall-clock bound. Unsupported inference remains unknown.
Exact finite-dimensional quantum lower witnesses are the next prerequisite
for general proved positive quantum-gap claims. Canonical transport, batch
execution, template reuse, and a searchable evidence ledger then support
repeated experiments without silently changing the mathematical question.

\section{Exact finite-dimensional quantum lower witnesses}
\label{sec:exactquantum}

A quantum upper certificate alone cannot establish an advantage over the local
value. A complementary certificate is an explicit admissible realization whose
score is evaluated exactly. We represent density matrices and local POVM effects
over $\mathbb Q+i\mathbb Q$, retaining the original local tensor dimensions.
Exact Hermiticity, state trace, positivity, and POVM completeness establish a
finite tensor-product strategy. Its rational Born score $s$ then proves
$\val_{\Q_{\rm fin}}(G)\geq s$ for a rational verifier.

\subsection{Complex positivity and tensor order}
For a Hermitian matrix $H=A+iB$, where $A^T=A$ and $B^T=-B$, define
\[
 \mathcal R(H)=\begin{pmatrix}A&-B\\B&A\end{pmatrix}.
\]
For $z=x+iy$, direct expansion gives
\[
 z^\dagger Hz=\begin{pmatrix}x^T&y^T\end{pmatrix}
 \mathcal R(H)\begin{pmatrix}x\\y\end{pmatrix}.
\]
Hence $H\succeq0$ if and only if $\mathcal R(H)\succeq0$, which can be checked
with the rational Schur procedure of Section~\ref{sec:xorcertificates}. This is
only a verification representation. It does not double the physical local
systems, and tensoring the separate realifications is not the original complex
tensor model. The implementation first computes complex tensor products in the
declared factor order and then evaluates the Born trace exactly. Its imaginary
part must vanish, probabilities must be nonnegative, and each conditional row
must sum to one.

\subsection{Rounding and repair with explicit scope}
The use of numerical discovery followed by rational verification is related to
exact reconstruction techniques for semidefinite certificates
\cite{peyrl2008}. Here the constraints are state normalization and POVM
normalization, and the final object is a feasible strategy rather than an SOS
upper certificate. A numerical candidate is converted explicitly; no floating
entry is accepted by the exact matrix decoder.

Let $\widehat\rho$ be a rational Hermitian approximation in dimension $D$.
Set
\[
 \rho_0=\widehat\rho+\frac{1-\Tr\widehat\rho}{D}\id.
\]
This has trace one even when coarse rounding made the approximate trace zero.
For a $k$-outcome POVM with rational Hermitian approximations $\widehat E_a$, set
\[
 F_a=\widehat E_a+\frac{\id-\sum_b\widehat E_b}{k}.
\]
Then $\sum_aF_a=\id$. These affine corrections need not preserve positivity.
For rational $t,u\in[0,1]$, define
\[
 \rho(t)=(1-t)\rho_0+t\id/D,\qquad
 E_a(u)=(1-u)F_a+u\id/k.
\]
The same $u$ is used for all effects of each POVM. At $t=u=1$ the state and
measurements are exactly feasible. Exact PSD tests and a bounded bisection can
retain a feasible endpoint closer to the uncorrected candidate. The final
strategy is fully rechecked, and its score is recomputed; no numerical objective
or eigenvalue becomes a proof merely through rounding.

This repair preserves physical dimensions. It does not in general preserve a
fixed state, projectivity, product-state form, PPT restrictions, or independent
sources. Those promises require additional checks or a different construction.
The initial certificate schema therefore declares only the unrestricted finite
tensor-product model. Dimension bounds can be checked directly against an exact
point, but unsupported resource promises must not be inherited from its numerical
origin.

\begin{proposition}[Density for strict score comparisons]
For fixed finite local dimensions and finite input/output alphabets, rational
complex states and POVMs satisfying the exact constraints are dense in the
ordinary tensor-product strategy space. Consequently any strict score advantage
attained by such a strategy is also attained by a sufficiently close rational
feasible strategy.
\end{proposition}
\begin{proof}
Mix a state and its POVMs by arbitrarily small amounts with the normalized
identities. The resulting matrices are strictly positive. Rational Hermitian
approximations, followed by the affine trace and completeness corrections above,
can be chosen arbitrarily close to them. Sufficiently small perturbations remain
positive because the positive-definite cone is open. The Born score is continuous
in the finitely many matrix entries, so a strict score comparison is preserved.
\end{proof}
This statement does not give a uniform search-success, dimension, or runtime
bound. Precision and repair budgets must be allowed to increase for arbitrary
approximation. It also does not guarantee rational attainment of every perfect
or optimal strategy. For example, a rational verifier and rational matrix data
have a rational score, whereas the optimal CHSH winning probability is irrational.
The symbolic CHSH certificate remains a separate exact algebraic path.

\subsection{Reference certificates and implementation boundary}
A Bell-state density matrix with rational entries, Alice's observables $(Z,X)$,
and Bob's observables $(cZ+sX,cZ-sX)$ with
\[
 c=803761/1136689,\qquad s=803760/1136689,\qquad c^2+s^2=1
\]
give the exact CHSH score
\[
 \tfrac12+\tfrac14(c+s)=\frac{3880899}{4546756}>\tfrac34.
\]
The GHZ X/Y strategy and the chosen magic-square strategy also admit rational
complex data and achieve score one. An optional Julia standard-library script
checks these reference constructions independently with unbounded rational
integers; Julia is not a dependency of the Python verifier.

Portable certificates bind the verifier fingerprint, physical dimensions, full
state/POVM data, and claimed score. Decoding verifies canonical rational strings
under explicit bit and allocation budgets. All integer components are normalized
to unbounded integers, including values originally held in fixed-width array
scalars. Budget exhaustion is distinct from a failed positivity test. Rechecking
uses exact arithmetic and supplies a lower bound only; it establishes neither
optimality nor self-testing. The claim layer can combine this lower witness with
an independently checked classical upper bound to prove a positive quantum gap.

\section{Structural properties and witnessed perfect play}
\label{sec:properties}

A property query must specify the mathematical predicate before selecting an
analysis method. Structural predicates inspect rational verifier data: acceptance
rewards, Boolean winning conditions, and dependence only on the parity of binary
answers. An XOR score can have constant or signed parity rewards; an XOR winning
predicate accepts exactly one parity at each selected question. These notions
apply at arbitrary party count and do not themselves license the two-party vector
SDP. The question domain is explicit: all stored inputs or only the positive-mass
sampling support.

Perfect play concerns the sampled game and requires acceptance rewards on that
support. It is an attainment statement, not merely the equality of a supremum
with one. The initial property fragment considers ordinary local and finite
tensor-product quantum strategies. More restrictive resources and self-testing
claims require additional rules rather than an inferred Boolean result.

\begin{proposition}[Sufficient property evidence]
For an acceptance game, a verified finite quantum strategy with score one proves
finite quantum perfectness, while a checked quantum upper bound below one
excludes it. A verified perfect local strategy also gives a perfect finite quantum
strategy by using one-dimensional systems and identity/zero effects. A perfect
finite quantum strategy together with a checked local upper bound below one proves
pseudotelepathy. Either a perfect local strategy or exclusion of quantum perfectness
refutes pseudotelepathy.
\end{proposition}
\begin{proof}
The point certificates establish admissible attained scores directly. An upper
bound below one excludes every score-one realization in its stated model.
Deterministic local responses have the indicated finite quantum embedding.
Pseudotelepathy is the conjunction of finite quantum perfectness and the absence
of a perfect local strategy, so the remaining conclusions follow.
\end{proof}
A nonsignalling score-one point is not a quantum point. A generic lower bound of
one is likewise not an attainment certificate; nonclosed quantum correlation sets
make the distinction substantive \cite{slofstra2019}. Numerical rank and optimizer
uniqueness do not repair this missing implication. The property engine therefore
records checked attainment as a separate fact and keeps its originating resource.

Every admitted attained score also supplies a lower value bound. This fact is
included in consistency checking even if a future producer omits a separate lower
bound field. For acceptance games, checked lower bounds above one, upper bounds
below zero, or attained scores outside $[0,1]$ are contradictory premises and
quarantine the assessment. Loose bounds outside that interval in the opposite
direction are not contradictions. Absence of a proof remains unknown; it is not
evidence that the property is false.

All property definitions bind to the exact verifier and support convention.
Structural counterexamples record locations and reasons that can be replayed on
the input data. Evidence-backed assessments recheck raw artifacts through the same
admission layer as value claims. Invalid, unsupported, exhausted, or conflicting
checks produce operational outcomes separate from mathematical verdicts. Receipts
are not standalone proof tokens, and no general strategy-discovery procedure is
launched implicitly by property assessment.

\section{Checked local relabellings}
\label{sec:relabellings}

Relabellings remove coordinate choices without changing the experimental
information. They are a restricted part of the broader theory of local
transformations, which also includes irreversible maps and alphabet-changing
liftings \cite{rosset2020transformations}. A relabelling must move a resource
contract as well as a score table whenever party roles are asymmetric.

\begin{definition}[Directed relabelling]
Let $\pi$ assign each new party $i$ to an old party $\pi(i)$. Let $\tau_i$
be a bijection from its new questions to the old questions, and let
$\sigma_{i,x'}$ be a bijection of answers for each new local question $x'$.
For new profiles $(x',a')$, define the old profiles by
\[
 x_{\pi(i)}=\tau_i(x'_i),\qquad
 a_{\pi(i)}=\sigma_{i,x'_i}(a'_i).
\]
The transported game and behavior satisfy
\[
 \mu'(x')=\mu(x),\qquad V'(a',x')=V(a,x),\qquad
 p'(a'\mid x')=p(a\mid x).
\]
\end{definition}

Both verifier tables are retained separately, even outside the sampling support.
The inverse answer map at old question $x$ is
$\sigma_{i,\tau_i^{-1}(x)}^{-1}$; inverting each answer permutation without
changing its question index would be incorrect. For a second relabelling
$(\kappa,\eta,\theta)$ applied after the first, composition is
\begin{align*}
 \pi_{\mathrm{comp}}(k)&=\pi(\kappa(k)),\\
 \tau^{\mathrm{comp}}_k(z)&=\tau_{\kappa(k)}(\eta_k(z)),\\
 \sigma^{\mathrm{comp}}_{k,z}(c)&=
 \sigma_{\kappa(k),\eta_k(z)}(\theta_{k,z}(c)).
\end{align*}

\begin{proposition}[Resource-aware score equivalence]
Transport local dimension bounds by $d'_i=d_{\pi(i)}$, PPT cuts by
$\pi^{-1}$, fixed-state tensor factors by $\pi$, and source endpoints by
$\pi^{-1}$ while preserving source and endpoint-position order. With these
transported contracts, relabelling gives a score-preserving bijection of the
implemented strategy models. Consequently the resource values agree, and
existence of an attaining strategy is preserved.
\end{proposition}
\begin{proof}
Change variables by the input and output bijections in the finite score sum.
Every term, including its sampling weight, is unchanged. Deterministic responses
transform as $f'_i(x')=\sigma_{i,x'}^{-1}(f_{\pi(i)}(\tau_i(x')))$.
For quantum strategies use
$M'_{i,a'\mid x'}=M_{\pi(i),\sigma_{i,x'}(a')\mid\tau_i(x')}$
and the corresponding tensor-factor permutation of the state. Born traces,
positivity, and POVM sums are preserved. Dimensions and partial-transpose
conditions follow the same factor renaming. Independent source laws or states
are unchanged, with only their recipient labels renamed. The inverse maps
recover every strategy. Commuting models can rename their observable families,
and the same finite coordinate permutation preserves nonsignalling constraints
and closures. Equality of suprema follows without assuming their attainment.
\end{proof}

This statement concerns mathematical resource sets. Numerical matrix validation
can fail after an exact entry permutation because eigensolver rounding changes
at a tolerance boundary. The implementation reports that failure at the original
tolerance; it neither repairs the state nor enlarges its acceptance threshold.
Party-coordinate permutations are not asserted to be local physical swap channels.

Evidence transport verifies a raw source certificate, constructs new witness
data, and rechecks a fresh certificate against the target verifier. Deterministic
and exact quantum witnesses use the formulas above. For bipartite XOR scores,
answer flips conjugate the Gram cost by a diagonal sign matrix, while party and
question relabellings permute its nodes. The dual diagonal follows the node
permutation; target PSD verification establishes the new upper certificate.
Local optimality artifacts retain independent enumeration replay. Nonsignalling
equality-dual and oriented spiral-copy rows require additional translation rules
and are not transported by changing hashes. Provenance receipts remain separate
from evidence admitted by the claim engine.

\bibliographystyle{plainurl}
\bibliography{bibliography}
\end{document}
